\subsection{Aplicación al PFC de gestión de pre-sustentaciones}

En el PFC del equipo BCCP desarrollamos un sistema de gestión de
documentos de pre-sustentaciones de trabajos de titulación. En la
arquitectura del proyecto, el servicio de Evaluación y actas necesita
consultar al de Solicitudes y documentos antes de habilitar la rúbrica
para el tribunal~\cite{pfcBccpComunicacion}.

\paragraph{gRPC para consultar una solicitud.}
En el esqueleto del PFC, Evaluación llama al método
\texttt{ConsultarSolicitud} con el identificador de la solicitud.
Solicitudes devuelve su estado, si el informe y el reporte de URKUND
están cargados y el porcentaje de coincidencia. Usamos gRPC porque
Evaluación necesita una respuesta síncrona y un contrato que fije los
tipos de esos campos. El cliente espera hasta cinco segundos y
comunica el estado de error si falla la consulta. En el flujo previsto,
sin esos datos confirmados, Evaluación no habilita la
rúbrica~\cite{pfcBccpRepositorio,pfcBccpComunicacion}.

\paragraph{Sockets TCP para un diagnóstico interno.}
Como extensión de la práctica, proponemos una prueba de conectividad
desde el entorno de Evaluación hacia el puerto \(50051\) de Solicitudes,
que es el puerto del esqueleto local~\cite{pfcBccpRepositorio}.
Un cliente de diagnóstico podría abrir un socket TCP con un tiempo de
espera y registrar si la conexión se estableció, fue rechazada o expiró.
Ese resultado ayudaría a distinguir un problema de conectividad de un
error devuelto por la llamada gRPC. Una conexión establecida indica que
el puerto acepta conexiones. Para comprobar que el trámite puede
consultarse correctamente, necesitamos ejecutar la llamada gRPC.
